\section{ماژول admm\_pruner (ADMM Pruner with Distillation and Dynamic Rho)}
\label{sec:module_admm_pruner}

\subsection{معرفی و انگیزه}
این مکانیزم بر مبنای مبانی نظری پیشرفته در ادبیات توجه و فشرده‌سازی \cite{kim2022learned} توسعه یافته است.

هرس کردن وزن‌ها \lr{(Weight Pruning)} به عنوان یکی از موثرترین تکنیک‌های سبک‌سازی شبکه‌های عصبی شناخته می‌شود؛ با این حال، حذف ناگهانی یا مقداری خالص وزن‌ها \lr{(Magnitude Pruning)} معمولاً به از دست رفتن اتصالات حیاتی و زوال ظرفیت زبانی مدل می‌انجامد. علاوه بر این، تنک‌سازی نامنظم \lr{(Unstructured Sparsity)} به دلیل عدم تطابق با خطوط لوله حافظه در شتاب‌دهنده‌های گرافیکی، شتاب محاسباتی ملموسی در استنتاج ایجاد نمی‌کند.

ماژول فشرده‌سازی تخصصی \lr{ADMMPruner} یک چارچوب بهینه‌سازی دقیق مبتنی بر الگوریتم جهت‌های متناوب لاگرانژی افزوده \lr{(ADMM)} را به همراه تقطیر بازسازی لایه‌ای \lr{(Layer-Wise Reconstruction Distillation)}، تخمین‌گر گرادیان مستقیم \lr{(Straight-Through Estimator - STE)} و الگوی تنک‌سازی ساختاریافته ۲:۴ معرفی می‌نماید. در این ساختار، متغیر تنک در تعاملی پویا با وزن‌های پیوسته به سمت آرایش ماتریسی هرس‌شده با کمترین اتلاف خروجی هدایت می‌شود و با به‌روزرسانی پویای ضریب جریمه $\rho$، پایداری همگرایی تضمین می‌گردد.

\subsection{مبانی نظری و فرمول‌بندی ریاضی}
مسئله هرس لایه‌ای به عنوان بهینه‌سازی مقید زیر فرمول‌بندی می‌شود:
\begin{equation}
\min_{\mathbf{W}, \mathbf{Z}} \frac{1}{2} \|\mathbf{W} - \mathbf{W}_0\|_F^2 + \lambda_{\text{distill}} \mathcal{L}_{\text{distill}}(\mathbf{W}; \mathbf{X}_{\text{calib}}) + g(\mathbf{Z}) \quad \text{s.t.} \quad \mathbf{W} - \mathbf{Z} = \mathbf{0}
\end{equation}
که در آن $\mathcal{L}_{\text{distill}} = \|\mathbf{X}_{\text{calib}} \mathbf{W}^\top - \mathbf{X}_{\text{calib}} \mathbf{W}_0^\top\|_F^2$ خطای تقطیر خروجی لایه بر روی فعال‌سازی‌های کالیبراسیون و $g(\mathbf{Z})$ تابع مشخصه قید تنک‌سازی (الگوی ساختاریافته ۲:۴ یا تنک‌سازی بلوکی) است. تابع لاگرانژی افزوده متناظر عبارت است از:
\begin{equation}
\mathcal{L}_{\rho}(\mathbf{W}, \mathbf{Z}, \mathbf{U}) = \frac{1}{2} \|\mathbf{W} - \mathbf{W}_0\|_F^2 + \lambda_{\text{distill}} \mathcal{L}_{\text{distill}}(\mathbf{W}) + \frac{\rho}{2} \|\mathbf{W} - \mathbf{Z} + \mathbf{U}\|_F^2 + g(\mathbf{Z})
\end{equation}
گام‌های الگوریتم در هر تکرار $k$ به شرح زیر اجرا می‌شوند:

\paragraph{۱. به‌روزرسانی لایه‌ای وزن‌ها ($\mathbf{W}$):}
بهینه‌سازی بر روی $\mathbf{W}$ ترکیبی از تقریب وزن‌های مبنا، تقطیر گرادیانی فعال‌سازی و انطباق با متغیر دوگان است:
\begin{equation}
\mathbf{W}^{(k+1)} = \left( (1 + \rho)\mathbf{I} + 2\lambda_{\text{distill}} \mathbf{X}^\top \mathbf{X} \right)^{-1} \left( \mathbf{W}_0 + 2\lambda_{\text{distill}} \mathbf{Y}_{\text{orig}}^\top \mathbf{X} + \rho (\mathbf{Z}^{(k)} - \mathbf{U}^{(k)}) \right)
\end{equation}

\paragraph{۲. تصویر بر فضای تنک با برجستگی افزوده ($\mathbf{Z}$):}
برخلاف هرس ساده اندازه، متغیر تنک $\mathbf{Z}$ با تصویر بر اساس برجستگی ترکیبی وزن و متغیر دوگان $|\mathbf{W}^{(k+1)} + \mathbf{U}^{(k)}|$ گزینش می‌شود:
\begin{equation}
\mathbf{Z}^{(k+1)} = \Pi_{\mathcal{S}_{2:4}}\left( \mathbf{W}^{(k+1)} + \mathbf{U}^{(k)} \right)
\end{equation}
در این مرحله، یک ماسک باینری $\mathbf{M}^{(k+1)} \in \{0, 1\}^{d_{\text{out}} \times d_{\text{in}}}$ تشکیل شده و مقادیر متناظر نگهداری می‌شوند.

\paragraph{۳. به‌روزرسانی متغیر دوگان و ضریب جریمه پویا ($\rho$):}
متغیر دوگان خطای برازش را انباشته کرده و ضریب $\rho$ به منظور شتاب‌بخشی به همگرایی با نرخ افزایشی $\gamma > 1$ مقیاس می‌شود:
\begin{align}
\mathbf{U}^{(k+1)} &= \mathbf{U}^{(k)} + \mathbf{W}^{(k+1)} - \mathbf{Z}^{(k+1)} \\
\rho^{(k+1)} &= \min\left(\rho_{\max}, \rho^{(k)} \cdot \gamma\right)
\end{align}

\subsection{معماری و پیاده‌سازی}
پیاده‌سازی ماژول در فایل \lr{Codes/framework/compression\_engine/admm\_pruner.py} با ساختار صنعتی و مولفه‌های کلیدی زیر همراه است:
\begin{itemize}
\item \textbf{لایه خطی ماسک‌دار \lr{ADMMMaskedLinear}:} جایگزینی مستقیم لایه‌های خطی استاندارد با یک لایه مجهز به بافر ماسک غیرپایدار و اعمال تخمین‌گر مستقیم \lr{ADMMSTEPrune} جهت انتشار روان گرادیان‌ها در پس‌انتشار:
\begin{equation}
\mathbf{y} = \operatorname{Linear}\left(\mathbf{x}, \mathbf{W} \odot \mathbf{M}, \mathbf{b}\right)
\end{equation}
\item \textbf{بازکالیبراسیون نرمال‌سازی لایه \lr{(LayerNorm Recalibration)}:} تنظیم مجدد پارامترهای میانگین و واریانس در لایه‌های \lr{LayerNorm} مدل پس از اعمال ماسک، جهت اصلاح شیفت توزیع فعال‌سازی ناشی از حذف اتصالات.
\item \textbf{سازگاری کامل با سخت‌افزارهای تنسور کُر \lr{NVIDIA Ampere}:} استفاده از چیدمان ۲:۴ امکان بهره‌گیری حداکثری از شتاب دوبرابری پایپ‌لاین‌های محاسباتی ماتریسی مدرن را فراهم می‌سازد.
\end{itemize}

\subsection{نتایج تجربی}
عملکرد اعتبارسنجی ماژول \lr{admm\_pruner} بر روی بنچ‌مارک رسمی پایان‌نامه در جدول~\ref{tab:benchmark_admm_pruner} منعکس گردیده است:

\begin{table}[H]
\centering
\small
\caption{کارنامه تجربی و ارزیابی بنچ‌مارک سخت‌افزاری ماژول \lr{admm\_pruner} روی مدل‌های پایه}
\label{tab:benchmark_admm_pruner}
\begin{tabular}{lccccc}
\toprule
\textbf{مدل پایه \lr{(Base Model)}} & \textbf{دقت \lr{(Accuracy)}} & \textbf{شاخص \lr{F1}} & \textbf{حافظه \lr{(Peak VRAM)}} & \textbf{توان \lr{(Throughput)}} & \textbf{وضعیت} \\
\midrule
\lr{DistilBERT} & \lr{84.29\%} & \lr{0.8603} & \lr{551.24 MB} & \lr{43,663.37 tok/s} & \lr{PASS} \\
\lr{RoBERTa} & \lr{80.62\%} & \lr{0.7966} & \lr{336.39 MB} & \lr{13,333.06 tok/s} & \lr{PASS} \\
\lr{GPT-2} & \lr{80.73\%} & \lr{0.8146} & \lr{414.93 MB} & \lr{15,380.14 tok/s} & \lr{PASS} \\
\lr{Qwen} & \lr{78.21\%} & \lr{0.8201} & \lr{1,879.76 MB} & \lr{1,102.95 tok/s} & \lr{PASS} \\
\bottomrule
\end{tabular}
\end{table}

\subsection{محدودیت‌ها و تحلیل شکست}
مشاهدات آزمایشگاهی حاصل از ارزیابی هرس \lr{ADMM} ابعاد تحلیلی زیر را مشخص می‌سازد:
\begin{itemize}
\item \textbf{نرخ قبولی ۱۰۰ درصدی در کلیه مدل‌ها:} روش \lr{ADMMPruner} بر روی هر چهار معماری ناهمگن با موفقیت به وضعیت قبولی \lr{PASS} با سطح دقت بین \lr{78.21\%} تا \lr{84.29\%} نائل آمد.
\item \textbf{نقش تقطیر بازسازی لایه‌ای در حفظ دقت:} مقایسه عملکرد با روش‌های هرس مرسوم نشان می‌دهد که حضور جمله تقطیر $\mathcal{L}_{\text{distill}}$ و بازکالیبراسیون \lr{LayerNorm} نقش حیاتی در جلوگیری از فروپاشی شاخص F1 (بالای ۰٫۷۹ در تمام مدل‌ها) ایفا کرده است.
\item \textbf{تحلیل حساسیت به گام‌های کالیبراسیون:} در لایه‌های حساس مانند پیش‌بینی کلمات انتهایی، در صورتی که تعداد گام‌های تقطیر لایه‌ای (\lr{distill\_steps}) ناکافی باشد یا داده‌های ورودی کالیبراسیون تنوع دامنه زبانی لازم را پوشش ندهند، هرس ۲:۴ می‌تواند به اعوجاج الگوهای نحوی ظریف بیانجامد. همچنین مصرف حافظه در طول کالیبراسیون به دلیل فعال‌سازی ذخیره تانسورهای موقت نسبت به هرس ایستا اندکی بالاتر است.
\end{itemize}
